% Русская версия RC4. Исходная структура, метки и формулы сохранены.
\title{\textbf{Контрактно-относительная теория преобразований:}\\
\large От сертифицированного открытия к замыканию метаполитик}
\author{Александр Странд (Aleksandr Strand)\\[-1pt]{\small\href{mailto:strandlex@gmail.com}{strandlex@gmail.com}}}
\date{\vspace{-11pt}Ноябрь 2026 г.}
\maketitle

\begin{abstract}
Представлена \textbf{контрактно-относительная теория преобразований (CRTT)} --- механизированная система бюджетного открытия и сертифицированного изменения исследовательских политик. Исходя из достаточности информации относительно контракта, мы разрабатываем конструктивный синтез исследований, сертифицированный перенос методов, возобновляемый поиск морфизмов, компиляцию препятствий и платное приобретение свидетельств, достаточных для решения. Эти конструкции обеспечивают механизм замены с сохранением продолжения, отделяющий семантическую корректность от экономического выигрыша и ликвидности.

Главный результат устанавливает композиционное замыкание конечных корректно типизированных расписаний обновления метаполитик при фиксированных исходной и целевой задачах, фиксированном поставщике свидетельств и финансируемом начальном продолжении. Исполнимая реализация сохраняет членство в каноническом результате, точную характеризацию невозможности, баланс ресурсов и предшествующую сертифицированную историю. Принятые обновления строго улучшают локальное значение полной модельной стоимости; кумулятивная монотонность гарантируется для трасс без отклонённых обновлений. Контрпримеры показывают, почему корректное обучение и локально выгодные замены не обязательно дают глобальный экономический выигрыш.

Каноническая реализация включает 21 исходный файл Lean и 542 аудируемые декларации в Lean~4.36.0, независимые конечные перекрёстные проверки и девять исполнимых свидетелей релиза. Гарантии ограничены объявленной конечной детерминированной моделью; произвольное самопереписывание, автономное порождение свидетельств и улучшение физической производительности не предполагаются.
\end{abstract}

\section{Введение и происхождение исследования}
\label{sec:intro}
\subsection{От сертифицированной функциональной замены к препятствиям CRTT}
Исследовательская программа CRTT возникла из задачи \emph{сертифицированной функциональной замены} (Certified Functional Replacement, CFR): когда один внутренний механизм сложной вычислительной или нейронной системы можно заменить другим, сохранив предписанное контрактом поведение? Эксперименты CFR мотивировали различение трёх возможных барьеров замены:
\begin{enumerate}
    \item \textbf{Операторный барьер:} в классе кандидатов нет допустимого оператора, реализующего требуемое поведение.
    \item \textbf{Репрезентационный барьер:} у слоя наблюдения может отсутствовать общее допустимое действие, даже если каждое его состояние разрешимо по отдельности.
    \item \textbf{Барьер выбора:} даже при достаточном наблюдении механизм выбора может не находить допустимый оператор. Эмпирические сбои калибровки мотивируют такое различение, но сами по себе не доказывают репрезентационную невозможность.
\end{enumerate}

Формализация этих барьеров в бюджетной системе с проверяемыми доказательствами выявила три новых класса проблем. Они не являются попарными эквивалентами барьеров CFR, а возникают при переходе от статической замены к работающим исследовательским системам: \emph{структурная непереносимость} между задачами, \emph{затраты на приобретение и поиск} методов и свидетельств, \emph{ограничения миграции и ликвидности} при замене действующего контроллера. Отсюда главный исследовательский вопрос CRTT:
\begin{quote}
\emph{При каких условиях исследовательская система может изменять собственные методы и управляющие политики, сохраняя семантическую корректность, конечные ресурсы и происхождение результатов, подтверждённое доказательствами?}
\end{quote}

\subsection{Архитектура как последовательность переходов}
CRTT отвечает на этот вопрос посредством шести качественных математических переходов (рис.~\ref{fig:arch-map}). Каждый слой разрешает конкретное ограничение предыдущего (табл.~\ref{tab:genesis}). В каноническом стеке стоимости, остатки и резервы --- натуральные учётные величины объявленной детерминированной модели, а не физическое время или память.

\begin{figure}[h]
\centering
\small
\begin{tikzpicture}[
    node distance=0.35cm and 0.45cm,
    box/.style={rectangle, draw=black!75, rounded corners=2pt, align=center,
                text width=3.9cm, minimum height=1.25cm, fill=black!3},
    arrow/.style={->, >=stealth, thick}
]
\node[box] (b1) {\textbf{Раздел 2}\\Контрактная достаточность\\\emph{Что нужно знать?}};
\node[box, right=of b1] (b2) {\textbf{Раздел 3}\\Конструктивное исследование\\\emph{Как получить знание?}};
\node[box, right=of b2] (b3) {\textbf{Раздел 4}\\Компиляция препятствий\\\emph{Как превратить сбой в правило?}};
\node[box, below=0.55cm of b1] (b4) {\textbf{Раздел 5}\\Метасвидетельства для решения\\\emph{Как выбирать методы?}};
\node[box, right=of b4] (b5) {\textbf{Раздел 6}\\Сертифицированная замена\\\emph{Как менять контроллер?}};
\node[box, right=of b5] (b6) {\textbf{Раздел 7}\\Композиционное замыкание\\\emph{Как сохранить гарантии?}};
\draw[arrow] (b1) -- (b2);
\draw[arrow] (b2) -- (b3);
\draw[arrow] (b3.south) -- ++(0,-0.25) -| (b4.north);
\draw[arrow] (b4) -- (b5);
\draw[arrow] (b5) -- (b6);
\end{tikzpicture}
\caption{Концептуальное развитие CRTT: стрелки означают переход между исследовательскими вопросами, а не обязательно импорт исходников.}
\label{fig:arch-map}
\end{figure}

\begin{table}[h]
\centering
\footnotesize
\begin{tabular}{@{}>{\raggedright\arraybackslash}p{3.6cm}>{\raggedright\arraybackslash}p{5.2cm}>{\raggedright\arraybackslash}p{4.6cm}@{}}
\toprule
\textbf{Установленный результат} & \textbf{Обнаруженная граница} & \textbf{Следующий объект исследования} \\
\midrule
Функциональная заменимость & Допустимого оператора может не быть & Препятствие \\
Достаточность наблюдения & Требуемая информация заранее неизвестна & Стратегия наблюдения \\
Оптимальная стратегия & Повторный синтез в новой задаче дорог & Перенос \\
Переносимость & Требуется найти морфизм переноса & Конструктивное открытие \\
Полное открытие & Исчерпывающий поиск слишком дорог & Сертифицированное отсечение \\
Обучение правилам & Корректное обучение может не окупаться & Метаконтроллер \\
Выбор контроллера & Точные стоимости методов известны частично & Метасвидетельства \\
Достаточные метасвидетельства & Замена может нарушить продолжение & Сертифицированная миграция \\
Корректная миграция & Повторные замены взаимодействуют в истории & Композиционные законы трасс \\
Замена контроллера & Сама политика выбора может быть недостаточной & Замыкание метаполитик \\
\bottomrule
\end{tabular}
\caption{Исследовательские вопросы, обусловленные границами предшествующих конструкций.}
\label{tab:genesis}
\end{table}

\subsection{Путеводитель и дополнительный архив}
Разделы~\ref{sec:foundations}--\ref{sec:meta-evidence} строят конструктивный и экономический аппарат для формулировки и доказательства главных результатов разделов~6 и~7. Полные исходники Lean~4.36.0 (21 канонический файл, 542 аудируемые декларации), исторические записи этапов и девять свидетелей исполнения CLI содержатся в согласованном архиве релиза (\texttt{CRTT-v1.0-FULL-ARCHITECTURE-RECONCILED.zip}), служащем формальным приложением.

\section{Контрактно-относительные основания}
\label{sec:foundations}
\subsection{Контракты, достаточность и критерий существования}
Зафиксируем пространство состояний $X$, непустое пространство действий $A \neq \varnothing$ и предикат контракта $V : X \times A \to \mathrm{Prop}$. Для каждого $x \in X$ определим множество допустимых действий $A_C(x) = \{ a \in A \mid V(x, a) \}$. Отображение наблюдения $\phi : X \to Z$ задаёт слой $\phi^{-1}(z) = \{ x \in X \mid \phi(x) = z \}$ над каждой наблюдаемой сигнатурой $z \in \phi(X)$. Тотальная политика $\pi : Z \to A$ \emph{допустима относительно контракта через $\phi$}, если $V(x, \pi(\phi(x)))$ выполняется для всех $x \in X$.

\begin{theorem}[Критерий контрактно-относительной достаточности (\texttt{CRTT.selector\_iff})]
\label{thm:sufficiency}
Пусть $A \neq \varnothing$. Отображение наблюдения $\phi : X \to Z$ допускает тотальную контрактно-допустимую политику $\pi : Z \to A$ тогда и только тогда, когда каждый реализуемый слой имеет общее допустимое действие:
\begin{equation}
\label{eq:sufficiency}
\exists \pi : Z \to A,\; \forall x \in X,\; V(x, \pi(\phi(x)))
\iff
\forall z \in \phi(X),\; \bigcap_{\phi(x) = z} A_C(x) \neq \varnothing.
\end{equation}
\end{theorem}
\begin{proof}
($\Rightarrow$) Если $\pi$ удовлетворяет $V(x, \pi(\phi(x)))$ для всех $x$, то для любого $z \in \phi(X)$, $\pi(z) \in \bigcap_{\phi(x)=z} A_C(x)$.
($\Leftarrow$) Для каждого $z \in \phi(X)$ выберем свидетель $a_z \in \bigcap_{\phi(x)=z} A_C(x)$ (в общем случае с использованием классического выбора). Для $z \notin \phi(X)$ положим $\pi(z) = a_0$, где $a_0 \in A$ существует в силу $A \neq \varnothing$. Полученная тотальная политика удовлетворяет $V(x, \pi(\phi(x)))$ для всех $x \in X$.
\end{proof}
Без $A \neq \varnothing$ критерий неверен при $X = \varnothing$ и $Z \neq \varnothing$. Кроме того, этот ранний общий результат использует классический выбор, в отличие от конструктивного конечного поиска канонического стека, аудируемые декларации которого не зависят от \texttt{Classical.choice}.

\subsection{Коллективные препятствия и грубейшие разбиения}
Из критерия~\eqref{eq:sufficiency} следует, что достаточность нарушается, если некоторый слой имеет пустое пересечение допустимых действий. Важно, что уточнение $\phi$ не всегда устраняет препятствие: если $A_C(x) = \varnothing$ для некоторого $x \in X$, даже полное наблюдение состояния не даёт допустимого действия.

Даже если каждое состояние разрешимо ($A_C(x) \neq \varnothing$), препятствия не обязательно попарны. Определим совместимость миров:
\[
x \sim_C y \iff A_C(x) \cap A_C(y) \neq \varnothing.
\]
Рассмотрим три мира с множествами допустимых действий $A_C(x_1)=\{a,b\}$, $A_C(x_2)=\{b,c\}$ и $A_C(x_3)=\{a,c\}$ (\texttt{SufficientPartitions.triangle\_compat\_all}, \texttt{triangle\_no\_common\_action}). Любые два мира имеют общее действие, но $\bigcap_{i=1}^3 A_C(x_i) = \varnothing$: препятствие строго коллективное. Это явление ограничивает существование наибольшего достаточного разбиения в порядке уточнения (а не кодировки минимального размера).

\begin{theorem}[Существование грубейшего достаточного разбиения]
\label{thm:coarsest}
Грубейшее достаточное разбиение $X$ существует тогда и только тогда, когда $\sim_C$ является отношением эквивалентности и разбиение, порождённое $\sim_C$, само достаточно.
\end{theorem}
\begin{proof}[Эскиз доказательства]
Если $\sim_C$ --- достаточное отношение эквивалентности, любое достаточное разбиение обязано разделять несовместимые миры и, значит, уточняет $\sim_C$; следовательно, $\sim_C$ является грубейшим. Обратно, пусть грубейшее достаточное разбиение $\mathcal{P}^*$ существует. Достаточность $\mathcal{P}^*$ влечёт $A_C(w)\ne\varnothing$ для каждого $w$. Тогда любую совместимую пару $x\sim_C y$ можно объединить, оставив остальные состояния одноэлементными блоками. Поскольку любое достаточное разбиение уточняет $\mathcal{P}^*$, $x$ и $y$ лежат в одном классе. С другой стороны, достаточность $\mathcal{P}^*$ требует совместимости любых двух состояний одного класса. Следовательно, $\sim_C$ совпадает с отношением эквивалентности $\mathcal{P}^*$ и достаточно.
\end{proof}
Соответствующая декларация --- \texttt{coarsest\_exists\_iff} в \texttt{SufficientPartitions.lean}. В примере с треугольником $\sim_C$ --- универсальное отношение (а значит, эквивалентность), но оно недостаточно; грубейшего достаточного разбиения нет. Напротив, для точной цели $g:X\to Y$ равенство целевых значений само является грубейшим достаточным отношением (\texttt{exact\_target\_coarsest}).

\newpage
\subsection{Экономические основания: накопление библиотеки и минимаксное сожаление}
Установив достаточность, исследуем, когда накопление методов или приобретение наблюдений окупается при конечных бюджетах:
\begin{itemize}
    \item \textbf{Идеализированное накопление и строгий выигрыш.} Для фиксированного предиката допустимой трассы, если повторное использование не дороже открытия и дополнительная плата за поиск отсутствует, бюджетная достижимость монотонна по включению библиотек. Обозначим через $C_{\mathrm{ind}}^*$ и $C_{\mathrm{cum}}^*$ достигаемые минимумы по одним и тем же допустимым планам, причём $C_{\mathrm{ind}}(p)=C_{\mathrm{cum}}(p)+S(p)$ для каждого допустимого $p$. Строгий кумулятивный выигрыш имеет место тогда и только тогда, когда некоторый план $p$ даёт сертифицированную экономию $S(p)$, превышающую разрыв до независимого оптимума:
    \begin{equation}
    \label{eq:t15}
    C_{\mathrm{cum}}^* < C_{\mathrm{ind}}^* \iff \exists p :\; S(p) > C_{\mathrm{ind}}(p) - C_{\mathrm{ind}}^*.
    \end{equation}
    \item \textbf{Обязательная плата за поиск и онлайн-пороги.} Обязательная плата за поиск может нарушать монотонность: расширение библиотеки увеличивает стоимость с $2$ до $3$ при бюджете $2$. В модели скрытого метода неблагоприятно выбранный повторяющийся скрытый метод даёт сожаление $D-R$ для детерминированного первого выбора, когда цена открытия равна $D$, повторного использования --- $R\le D$ и имеется не менее двух задач. В отдельной пороговой модели повторяющейся задачи с ценой аренды $L$ и $R<L\le D$ стратегия точки окупаемости $2$-конкурентна относительно офлайн-оптимума.
    \item \textbf{Минимаксное сожаление и ценность наблюдения.} Для конечного списка сценариев, непустого конечного списка действий и разрешимого равенства наблюдений пусть $c(a,x)\in\mathbb N$. Определим сожаление действия $r(a,x) = c(a,x) - \min_{b \in A} c(b,x)$ и минимаксное сожаление наблюдения $\phi$ в худшем случае:
    \begin{equation}
    \label{eq:t18-regret}
    R^*(\phi) = \max_{z \in \phi(X)} \min_{a \in A} \max_{x : \phi(x)=z} r(a,x).
    \end{equation}
    Максимум пустого списка сценариев равен нулю. Минимаксный критерий характеризует $\varepsilon$-достаточность существованием в каждом реализуемом слое общего действия с сожалением не более $\varepsilon$. В примере XOR отсутствие наблюдения или наблюдение любого одного бита даёт сожаление $1$, обоих битов --- $0$. Приобретение фиксированного уточнения $\phi'$ за цену $c_{\mathrm{obs}}$ строго выгодно в точности тогда, когда
    \begin{equation}
    \label{eq:t18-acq}
    c_{\mathrm{obs}} + R^*(\phi') < R^*(\phi).
    \end{equation}
\end{itemize}

\subsection{От достаточности к исследованию}
Эти результаты включают онлайн-пороговую модель, но ещё не дают общего синтезатора деревьев ограниченной глубины и типизированного механизма переноса задач, разрабатываемых ниже. Это мотивирует слой конструктивного исследования и переноса (раздел~3).

Результаты оснований имеют отдельные исторические доказательные артефакты. Самодостаточный канонический стек устанавливает теорему релиза независимо от этих ветвей; их происхождение описано в разделе~8.
